% ============================================================================
% ANEXOS
% ============================================================================
% Los anexos contienen material complementario relevante para el trabajo pero
% que, por su extensión o nivel de detalle, no se incluye en el cuerpo principal.
% Ejemplos: código fuente extenso, manuales de instalación, encuestas, etc.
% ============================================================================

\chapter*{ANEXOS}
\addcontentsline{toc}{chapter}{ANEXOS}

% Se activa el formateo de anexos de LaTeX (cambia la numeración de Capítulos a letras A, B, etc.)
\appendix

\chapter{Manual de Instalación y Configuración}

[Describe en este anexo el paso a paso para instalar, configurar y ejecutar tu solución tecnológica...]

\section{Requisitos del sistema}
[Especifica los requisitos mínimos y recomendados de hardware, sistema operativo, bases de datos y entornos de ejecución.]

\section{Instrucciones de instalación y despliegue}
[Describe detalladamente los comandos, configuraciones de variables de entorno y pasos necesarios para desplegar la solución.]

\chapter{Código Fuente y Configuración Relevante}

[Puedes incluir listados de scripts de base de datos (SQL, NoSQL), archivos de configuración (Docker, variables de entorno, webpack, etc.) o fragmentos de código fuente que complementen el cuerpo de la obra.]
